\subsection{逆李代数的包络代数}

设 \(g\) 是 \(K\) 上的李代数，\(g^0\) 是逆李代数，\(\sigma\) 和 \(\sigma_0\) 分别是将 \(g\) 和 \(g^0\) 映到其包络代数 \(U\) 和 \(V\) 的典范映射。则 \(\sigma\) 是一个 \(\alpha\)-映射，从 \(g^0\) 到结合代数 \(U^0\)，即结合代数 \(U\) 的逆代数。因此存在唯一的同态 \(\phi\)，它从 \(V\) 到 \(U^0\)，将 1 映到 1，并满足 \(\sigma = \phi \circ \sigma_0\)。

命题 4. 同态 \(\phi\) 是从 \(V\) 到 \(U^0\) 的同构。

存在同态 \(\phi'\)，它将 U 映到 \(V^0\)，将 1 映到 1，并满足 \(\sigma_0 = \phi' \circ \sigma\)。可将 \(\phi'\) 视为从 \(U^0\) 到 V 的同态。于是 \(\sigma_0 = \phi' \circ \phi \circ \sigma_0\)，且 \(\sigma = \phi \circ \phi' \circ \sigma\)，从而 \(\phi' \circ \phi\) 和 \(\phi \circ \phi'\) 分别是 V 和 U 的恒等映射。故命题得证。

V 与 \(U^{0}\) 相同，这一识别由同构 \(\phi\) 给出。于是 \(\sigma_{0}\) 与 \(\sigma\) 相同。

在此同一视下，同构 \(\theta: x \mapsto -x\) 从 \(\mathfrak{g}\) 到 \(\mathfrak{g}^0\)，定义了同构 \(\tilde{\theta}\) 从 U 到 V = U \(^0\)。此同构可视为 U 的反自同构，称为 U 的主反自同构。若 \(x_{1}, \ldots, x_{n}\) 属于 \(\mathfrak{g}\)，则：

\[
\begin{array}{r l} \tilde {\theta} (\sigma (x _ {1}) \dots \sigma (x _ {n})) & = \tilde {\theta} (\sigma (x _ {n})) \dots \tilde {\theta} (\sigma (x _ {1})) = (- \sigma (x _ {n})) \dots (- \sigma (x _ {1})) \\ & = (- 1) ^ {n} \sigma (x _ {n}) \dots \sigma (x _ {1}). \end{array}\tag{2}
\]

